\scanpage{54}
\mspage{IV}

(i) Pour tout $X$ de t. f. sur $S$, $H^1(X,\unsure{G|C})$ est fini.

(ii) Pour tt sous-schéma \deleted{\unsure{fermé}} normal \deleted{de $S$} \inserted{intègre affine} $S_0$ de $S$, \deleted{\unsure{et}} tt \unsure{rev.} fini étale $S'$ de $S_0$, tout schéma \inserted{affine} lisse de t. fini $\unsure{X|X'}$ sur $S'$, à fibres géom. connexes non vides, $H^1(\deleted{\unsure{S}}X,G)$ est fini

(iii) Pour tt $S',X$ comme dessus, $S'\ne\varnothing$, il existe un ouvert non vide $U$ de $S'$ \unsure{tel|dt|d'où} que $H^1(X\mid U',G)$ soit fini.

Il \uuline{suffit} pour ceci qu'on ait ceci :

(iv) a) Pour tout $s\in S$, et tt schéma $\bar X$ de type fini lisse affine connexe sur $\overline{k(s)}$, $H^1(\bar X,G)$ est \unsure{fini|finie}

b) Pour tout \deleted{\unsure{$S$}} \deleted{\unsure{sous}} schéma affine normal intègre $S'$ \deleted{de $S$}, \inserted{\unsure{comme|nommé} dans ii)} et tt sous \unsure{gpe|groupe} $H$ de $G\times G$, \deleted{\illegible{1}} $H^1(S',H)$ est fini.

Lorsque $G$ est commutatif, on peut affaiblir b) en y remplaçant $H$ par $G$.

\margin{\rule{4em}{.4pt}}
\uline{Dém}. Equivalence de (i) à (iii) résulte aussitôt de lemme 3. Pour prouver qu'ils sont entraînés par (iv), on doit utiliser de plus lemme 1.
% archive sheet number: 53
